%sec:j
%1st paragraph
%1
\revised{Processing followed Section~\ref{sec:method}: two vertices defined each person region for trimming; Open3D extracted normals, dimensionality features, proportion of variance, and FPFH \cite{zhou2018open3d}.}
  %1-1 点座標のみは評価しない
  \revised{We do not evaluate point coordinates only on the basis of the comparison in existing research \cite{DBLP:journals/corr/QiYSG17}.}
    %1-1-1
    \revised{The research evaluated point coordinates only and point coordinates with point features.}
    %1-1-2
    \revised{The research reported higher accuracy with point coordinates and point features than with point coordinates only.}
\begin{deletedblock}
In our experiment, the proposed system was processed according to the methodology described in Section~\ref{sec:method}.
  Two vertices of the region where a person is located were determined and then the point clouds were trimmed. 
The features were generated using the process explained in Section~\ref{sec:method}. 
  Specifically, the normal estimator from Open3D was used and normals were stored with the coordinate data in point cloud data files. 
  The covariances were generated and then the dimensionality features and the proportion of variance were calculated. 
  The FPFH estimator from Open3D was utilized and the FPFH values were stored in a separate file from the coordinate data. 
\end{deletedblock}
  %1-2
  \revised{DBSCAN and k-means removed outliers as in Section~\ref{sec:method}; point counts were adjusted to 4,096 after clustering and downsampling \cite{pytorchpointnet++}.}
  \begin{deletedblock}
  The data were clustered using the process explained in Section~\ref{sec:method}. 
    Specifically, outliers and some clusters were removed after implementing DBSCAN.
    Some clusters were also removed after performing k-means. 
  After the clustering and downsampling, the number of points in each point cloud was adjusted to 4,096, which is the default value of the Github repository \cite{pytorchpointnet++}. 
  \end{deletedblock}
  %1-3
  Normalization was performed using different values for each point cloud dataset used for the training. 
  %1-4
  \deleted{PointCutMix augmented normalized training data by replacing the top $\lambda$\% of nearest-neighbor-matched points between three point clouds A, B, and C, with $\lambda$ drawn from a beta distribution \cite{ZHANG202258}.}
  \revised{PointCutMix augmented normalized training data by replacing the top $\lambda$\% of nearest-neighbor-matched points among three point clouds, with $\lambda$ drawn from a beta distribution \cite{ZHANG202258}.}
  %1-4-1
  \deleted{PointCutMix increases shape diversity among training samples and improves robustness to subject-dependent point-cloud density.}
  \revised{PointCutMix increases shape diversity among training samples and improves robustness to experiment-participant-dependent point cloud density.}
  \begin{deletedblock}
  Data augmentation was performed on the normalized data. 
    For data augmentation, PointCutmix was used\cite{ZHANG202258}. 
        Three point clouds A, B, and C were prepared. 
        The beta function was used to randomly determine a real number $\lambda$. 
        Points from point cloud B were assigned to each point in point cloud A. 
            For each point in A, the distance to all points in B was measured and the index of the closest point in B was assigned. 
            Then, the top $\lambda$\% of points in A, sorted by distance, were replaced with their assigned points from B. 
        Points from point cloud C were assigned to each point in point cloud A. 
            For each point in A, the distance to all points in B was measured and the index of the closest point in C was assigned. 
            Then, the top $\lambda$\% of points in A, sorted by distance, were replaced with their assigned points from C. 
  \end{deletedblock}

%2nd paragraph
%2
\deleted{Data were split 4:1 for training and validation; multi-scale grouping (MSG) formed PointNet++ groups \cite{pytorchpointnet++,PytorchPointnetPointnet2}.}
\begin{revisedblock}
We split frames of each of the three experiment participants at a ratio of 4:1.
  %2-1
  Four fifths of the frames were used for training.
  %2-2
  One fifth of the frames were used for validation.
%3
Multi-scale grouping (MSG) formed PointNet++ groups \cite{pytorchpointnet++,PytorchPointnetPointnet2}.
\end{revisedblock}
\begin{deletedblock}
The remaining data were split into 4:1 for the training and validation. 
Multi-scale grouping (MSG) was used to create groups in PointNet++, as based on prior methods \cite{pytorchpointnet++,PytorchPointnetPointnet2}.
An output containing the parameters of the classification was created after the training was completed. 
\end{deletedblock}
